\section{Related Work}
\label{sec:related_work}

\subsection{UAV-based Cross-View Geo-Localization}

UAV-based retrieval requires discriminative representations despite large changes in viewpoint, scale, and appearance. LPN~\cite{wang2022lpn} uses square-ring partitions to capture local patterns and surrounding context, while FSRA~\cite{dai2022fsra} aligns regions obtained from Transformer features across views. Sample4Geo~\cite{deuser2023sample4geo} improves representation learning through symmetric contrastive learning and hard negative sampling.

Generalization to unseen geographic regions is another important consideration. DAC~\cite{xia2024dac} combines cross-view feature alignment with scene consistency, and InfoGeo~\cite{zhang2026infogeo} aligns object-centric structural relations while suppressing view-specific information. Pretrained visual models also provide transferable representations for retrieval. In the related task of visual place recognition, AnyLoc~\cite{keetha2023anyloc} aggregates self-supervised features, while SALAD~\cite{izquierdo2024salad} fine-tunes DINOv2 and aggregates local features through optimal transport. These studies motivate the use of pretrained visual representations in UAV-satellite retrieval.

For efficient deployment, MEAN~\cite{chen2025mean} uses a lightweight backbone with multi-level feature enhancement and cross-view alignment, while SpikeViMFormer~\cite{chen2026spikevimformer} adopts a spike-driven Transformer backbone. Knowledge distillation provides another route to compact models. DepthGeoLoc~\cite{sun2026depthgeoloc} transfers depth-aware structural knowledge through coarse-to-fine distillation. MobileGeo~\cite{sun2025mobilegeo} uses hierarchical distillation and uncertainty-aware prediction alignment, together with view selection at inference. EfficientGeo also uses distillation to train a lightweight model, aiming for accurate retrieval in unseen geographic regions under onboard computational constraints.

\subsection{Knowledge Distillation}

A large capacity gap can make it difficult for a compact student to learn directly from a stronger teacher~\cite{cho2019efficacy}. Teacher assistant distillation~\cite{mirzadeh2020takd} introduces an intermediate-sized model to bridge this gap, while densely guided knowledge distillation~\cite{son2021dgkd} extends this strategy to multiple assistants. Inspired by these studies, EfficientGeo uses a Transformer-based intermediate teacher to transfer task-adapted foundation knowledge to a lightweight convolutional student.

Knowledge can also be transferred through relationships among samples. DarkRank~\cite{chen2018darkrank} and RKD~\cite{park2019rkd} distill similarity-based rankings and geometric relationships, respectively. For retrieval, Margin-MSE~\cite{hofstatter2020marginmse} matches score differences between relevant and non-relevant passages. GeoMFD~\cite{chen2026geomfd} transfers positive-negative similarity margins on teacher-selected hard negatives in both retrieval directions to reduce forgetting during continual adaptation. In our setting, the task-adapted foundation teacher transfers positive-negative similarity margins to a smaller intermediate teacher, using hard negatives selected by the foundation teacher.

Transferring features across different architectures is also challenging because Transformer and convolutional models represent visual information differently. Existing methods use projection modules~\cite{liu2022crossarchitecture,hao2023ofa} or combine heterogeneous teacher and student modules~\cite{li2025fbt} to facilitate transfer. VkD~\cite{miles2024vkd} aligns teacher and student features using an orthogonal projection and task-specific normalization. In EfficientGeo, the convolutional student learns compact descriptor targets obtained from two fixed projections of the intermediate teacher. The teacher subspace captures principal directions of its representations, while an independently sampled random subspace provides a second target. This design does not require correspondence between intermediate layers of the Transformer teacher and the convolutional student.
